%
% spec-audit-durability.tex
%
% Audit Line Durability: A Z Specification
% A formal model of the life of a sealed audit line across the actor set
% that can touch a checkout's live audit zone: a commit in the writing
% checkout, a commit in a sibling worktree, `git worktree remove`, session
% end, `ethos disable`, and `ethos session purge`. Current as of the code
% cited throughout (internal/hook/audit_seal.go, audit_vacuum.go,
% session_end.go; internal/session/store.go; internal/enable/enable.go).
%
% Compile: pdflatex spec-audit-durability.tex
%
% Must be .tex, not .md -- see spec-mission-lifecycle.tex's own header note:
% /z-spec:test and /z-spec:model_check glob for .tex specifically.
%

\documentclass[a4paper,10pt,fleqn]{article}
\usepackage[margin=1in]{geometry}
\usepackage{fuzz}
\usepackage{booktabs}
\usepackage{tabularx}
\usepackage[colorlinks=true,linkcolor=blue,citecolor=blue,urlcolor=blue]{hyperref}

\newcommand{\fpath}[1]{\texttt{#1}}
\newcommand{\field}[1]{\texttt{#1}}

\begin{document}

\title{Audit Line Durability: A Z Specification}
\author{Mike S, for Punt Labs}
\date{September 2026}
\hypersetup{
  pdftitle={Audit Line Durability: A Z Specification},
  pdfauthor={Punt Labs},
  pdfsubject={Z Specification},
  pdfcreator={fuzz/probcli}
}
\maketitle
\tableofcontents
\newpage

% ===================================================================
\section{Scope}
\label{sec:scope}
% ===================================================================

Every audit line an agent writes begins in a gitignored, per-checkout
live file, \fpath{<repo>/.punt-labs/local/ethos/sessions/<session-id>.audit.jsonl}
(docs/audit-seal.md \S Two zones, lines 65--73). \field{SealRepo}
(internal/hook/audit\_seal.go:49) moves a checkout's pending tail into a
tracked chunk file, and only a git commit in \emph{that} checkout makes the
chunk durable in history. Between those two events the line's only home is a
plain file on a filesystem that git does not protect for gitignored content,
and \fpath{git worktree remove} deletes exactly that filesystem without
consulting ethos at all: no hook in this repository fires on worktree
removal (a search of \fpath{internal/} confirms \fpath{internal/resolve}
only \emph{reads} worktree layout, at resolve.go:757--827, to find which
checkout owns a store; nothing reacts to one being torn down).

This specification models one audit line's journey through three states
--- spooled in the live zone, sealed into an untracked-but-tracked-path
chunk, committed into git history --- and checks a single safety property
against the actor set the mission brief names: a commit in the writing
checkout, a commit in a sibling worktree, \fpath{git worktree remove},
session end, \fpath{ethos disable}, and \fpath{ethos session purge}. It is a
baseline, not a design proposal: \S\ref{sec:state}--\ref{sec:ops} model the
system exactly as \fpath{internal/hook} and \fpath{internal/session}
implement it today. \S\ref{sec:alternatives} then evaluates three amendments
against the same property, each as an independent model variant.

\subsection{Modelling choice: one line, one writer checkout}

The real live file accumulates many lines over a session's life, and a repo
may hold arbitrarily many worktrees. Neither multiplicity changes the shape
of the property: \field{SealRepo} drains a checkout's \emph{entire} pending
tail in one atomic write (audit\_seal.go:201--219, ``\emph{tail selection...
write it whole to a new chunk}''), never line-by-line, so the state that
matters is an aggregate count, not a sequence; and every actor in the brief
acts on \emph{one} specific checkout at a time (the committing one, or the
one a roster names as the writer) regardless of how many other checkouts
exist. This licenses the model: \field{AuditLine} (\S\ref{sec:state})
tracks one small pool of lines belonging to one writer checkout, and a
second, unnamed checkout stands in for however many siblings the repo
actually has --- \field{CommitInSiblingWorktree}'s effect does not depend on
\emph{how many} siblings commit, only on whether one does.

\subsection{Non-goals}
\label{sec:nongoals}

\begin{itemize}
  \item \textbf{Per-line ordering and per-line watermarks.} \field{audit.SelectLiveTail}
    and \field{audit.Watermark} (audit\_seal.go:194--201) resolve exactly
    which lines are new since the last seal; this specification does not
    re-derive that logic, it assumes the aggregate count SealRepo would
    seal is known and models only what happens to that count.
  \item \textbf{Chunk corruption and quarantine.} \field{ReadChunkVerified}'s
    per-chunk hash check (audit\_seal.go:189--193) and the quarantine escape
    it feeds are a data-integrity concern once a line is already sealed;
    orthogonal to whether a line survives at all.
  \item \textbf{Counting worktrees.} Following the modelling choice above, a
    boolean \field{writerExists} stands for ``the writer checkout is still
    present''; the model does not carry a count of how many other worktrees
    exist, only that a sibling-worktree commit can happen.
  \item \textbf{git's own dirty-worktree refusal, except where a design
    variant depends on it.} Plain \fpath{git worktree remove} (no
    \texttt{--force}) refuses when the worktree has modified or untracked
    tracked-path content, but not when the only changes are to gitignored
    paths --- which is exactly what the live zone is (docs/audit-seal.md
    \S Two zones: \fpath{.punt-labs/local/} is gitignored). The current
    design (\S\ref{sec:state}--\ref{sec:ops}) never leaves a
    \emph{tracked}, uncommitted chunk sitting in the writer checkout long
    enough for this to matter, because sealing only ever happens
    immediately before a commit in the same checkout
    (\S\ref{sec:opcommitwriting}); it becomes load-bearing only in Design~3
    (\S\ref{sec:design3}), where it is modelled explicitly.
  \item \textbf{Re-enabling ethos.} \fpath{ethos enable} after a disable is
    the inverse of \field{EthosDisable} (\S\ref{sec:opdisable}) and is not
    one of the six named actors; omitted for economy.
\end{itemize}

% ===================================================================
\section{Free Types}
\label{sec:freetypes}
% ===================================================================

\subsection{ZBOOL}

fuzz has no built-in boolean, and \field{BOOL}/\field{true}/\field{false}
collide with B keywords under ProB's translation:

\begin{zed}
ZBOOL ::= ztrue | zfalse
\end{zed}

Four independent booleans ride this type in \field{AuditLine}
(\S\ref{sec:state}): whether the writer checkout still exists on disk,
whether its pre-commit seal hook is installed, whether some durable
record (a live roster entry or a flagged purge tombstone) still names the
session, and whether an unwarned, unsealed loss has ever occurred on this
line's trajectory.

% ===================================================================
\section{Global Constants}
\label{sec:constants}
% ===================================================================

The pool of lines tracked is bounded to keep ProB's state space finite, per
the z-spec ProB-compatibility guidance on bounding numeric variables. The
shape of every predicate below is independent of the concrete value.

\begin{axdef}
totalLines : \nat
\where
totalLines = 2
\end{axdef}

% ===================================================================
\section{State}
\label{sec:state}
% ===================================================================

\begin{schema}{AuditLine}
spooled, sealed, committed : \nat \\
writerExists, hookInstalled, rosterKnown, lostUnwarned : ZBOOL
\where
spooled + sealed + committed \leq totalLines \\
writerExists = zfalse \implies (spooled = 0 \land sealed = 0)
\end{schema}

Two predicates:

\begin{itemize}
  \item \field{spooled + sealed + committed} $\leq$ \field{totalLines} is a
    conservation bound, not a code invariant: it fixes a small pool of
    lines and lets their location --- spool, chunk, or history --- account
    for at most that many at once. The \emph{deficit},
    \field{totalLines} $-$ (\field{spooled} $+$ \field{sealed} $+$
    \field{committed}), is how many lines this trajectory has lost; there
    is deliberately no separate counter for it; see
    \S\ref{sec:opworktreeremove}.
  \item \field{writerExists = zfalse} $\implies$
    (\field{spooled = 0} $\land$ \field{sealed = 0}) is definitional, not
    the safety property: it says what ``the writer checkout is gone''
    \emph{means} for a live-zone file and an unstaged-or-staged chunk file
    that both live only inside that checkout's working tree. It is not
    baked into any operation's guard; every operation that clears
    \field{writerExists} must independently clear \field{spooled} and
    \field{sealed} to stay well-typed, and \S\ref{sec:opworktreeremove}
    shows the real operation does so by \emph{destroying} them, not by
    having sealed or reported them first.
\end{itemize}

\field{lostUnwarned} carries no invariant of its own. This is deliberate,
and mirrors an errata already on record in spec-mission-lifecycle.tex
\S Validation: a state invariant of the form
$\field{lostUnwarned} = \field{zfalse}$ living inside \field{AuditLine}'s
own \texttt{where}-clause would, under Z's schema calculus, make any
$\Delta \field{AuditLine}$ operation that tries to set
$\field{lostUnwarned}' = \field{ztrue}$ \emph{unsatisfiable} rather than
\emph{caught} --- ProB would report no violation for exactly the reason
there is nothing left to explore, which is a false negative, not a clean
bill of health. \field{lostUnwarned} is instead a plain tracked witness:
the safety property (\S\ref{sec:safety}) is that no reachable state has
\field{lostUnwarned = ztrue}, checked by inspecting the state space
ProB actually explores, not by an invariant that could silently swallow
the case it exists to catch.

% ===================================================================
\section{Initialization}
\label{sec:init}
% ===================================================================

\field{Init} starts with the whole small pool already sitting unsealed in
the live zone of a fresh, present writer checkout with its seal hook
installed and its session still known to the roster --- the state a session
is in for most of its life, and the state every actor in \S\ref{sec:ops}
must handle safely.

\begin{schema}{Init}
AuditLine'
\where
spooled' = totalLines \\
sealed' = 0 \\
committed' = 0 \\
writerExists' = ztrue \\
hookInstalled' = ztrue \\
rosterKnown' = ztrue \\
lostUnwarned' = zfalse
\end{schema}

% ===================================================================
\section{Operations}
\label{sec:ops}
% ===================================================================

\subsection{CommitInWritingCheckout}
\label{sec:opcommitwriting}

Models \field{SealRepo} (audit\_seal.go:49--64) running at pre-commit in
the checkout that holds the line, immediately followed by the commit
landing (docs/audit-seal.md: ``\emph{the freshly staged chunks land in the
same commit as the work they document}''). Requires the hook still be
installed: with no pre-commit hook there is nothing to invoke
\field{SealRepo} at all (\S\ref{sec:opdisable}).

\begin{schema}{CommitInWritingCheckout}
\Delta AuditLine
\where
writerExists = ztrue \\
hookInstalled = ztrue \\
writerExists' = writerExists \\
hookInstalled' = hookInstalled \\
rosterKnown' = rosterKnown \\
spooled' = 0 \\
sealed' = 0 \\
committed' = committed + spooled + sealed \\
lostUnwarned' = lostUnwarned
\end{schema}

\subsection{CommitInWritingCheckoutNoHook}

The companion case: a commit lands in the writer checkout, but its
pre-commit hook was removed by an earlier \field{EthosDisable}
(enable.go:193, \field{unchainHooks}), so \field{SealRepo} is never
invoked. The commit happens; nothing about the audit line does.

\begin{schema}{CommitInWritingCheckoutNoHook}
\Xi AuditLine
\where
writerExists = ztrue \\
hookInstalled = zfalse
\end{schema}

\subsection{CommitInSiblingWorktree}

Models a commit in a \emph{different} checkout of the same repository.
\field{SealRepo} is repo-wide only in the sense of visiting every session
and mission directory \emph{inside the committing checkout}
(\field{listRepoSessions(repoRoot)}, audit\_seal.go:68, where
\field{repoRoot} is always the committing checkout ---
audit\_seal.go:49--50). It has no way to reach another checkout's live
zone, so this operation cannot touch \field{spooled} or \field{sealed} at
all. It does run \field{VacuumCrossCheck} (audit\_vacuum.go:46), which
warns on stderr when it can see the writer checkout is gone
(audit\_vacuum.go:96--107) --- but that is a message emitted \emph{after}
any loss already happened via a separate transition
(\S\ref{sec:opworktreeremove}), not a state change, and not something that
can retroactively make the destroying transition itself compliant with
``sealed or warned \emph{first}''. It is therefore modelled as a pure
query, changing nothing:

\begin{schema}{CommitInSiblingWorktree}
\Xi AuditLine
\end{schema}

\subsection{WorktreeRemove}
\label{sec:opworktreeremove}

Models \fpath{git worktree remove}. As \S\ref{sec:scope} establishes, no
ethos hook fires on this event; the guard is only git's own precondition
that there is a checkout there to remove. The operation deletes the
checkout's working tree wholesale, so whatever was in the live zone or an
uncommitted chunk file goes with it. \field{lostUnwarned} latches to
\field{ztrue} exactly when this destroys a nonempty pool, and never
resets --- there is no operation anywhere in this specification that turns
\field{ztrue} back to \field{zfalse}, matching the fact that a destroyed
line cannot un-destroy itself no matter what happens afterward.

\begin{schema}{WorktreeRemove}
\Delta AuditLine
\where
writerExists = ztrue \\
writerExists' = zfalse \\
spooled' = 0 \\
sealed' = 0 \\
committed' = committed \\
hookInstalled' = hookInstalled \\
rosterKnown' = rosterKnown \\
(spooled > 0 \lor sealed > 0) \implies lostUnwarned' = ztrue \\
(spooled = 0 \land sealed = 0) \implies lostUnwarned' = lostUnwarned
\end{schema}

\subsection{SessionEnd}

Models \field{HandleSessionEnd} (session\_end.go:46--67): it clears the
session's roster entry via \field{ss.Delete} (session\_end.go:57), which in
turn only removes the roster file and mission sidecars
(internal/session/store.go:278--301, \field{deleteFiles}) --- no seal, and,
per \field{HandleSessionEnd}'s own doc comment, no unsealed-line probe,
deliberately: ``\emph{failing the hook here would not improve on that
guarantee... and would only make session end itself unreliable}''
(session\_end.go:41--45). The live audit file and any uncommitted chunk are
untouched; only the bookkeeping that a \emph{later} check could use to
notice their loss disappears.

\begin{schema}{SessionEnd}
\Delta AuditLine
\where
spooled' = spooled \\
sealed' = sealed \\
committed' = committed \\
writerExists' = writerExists \\
hookInstalled' = hookInstalled \\
rosterKnown' = zfalse \\
lostUnwarned' = lostUnwarned
\end{schema}

\subsection{EthosDisable}
\label{sec:opdisable}

Models \fpath{ethos disable} (internal/enable/enable.go): it removes the
pre-commit hook chain (\field{unchainHooks}, enable.go:193) and, if any
lines are outstanding, prints an informational --- not a refusal ---
line (enable.go:197--198: ``\emph{N unsealed audit lines remain in the
local zone; re-enable to seal them}''). Nothing about the pool itself
changes; only the hook that would otherwise drain it on the next commit in
this checkout is gone.

\begin{schema}{EthosDisable}
\Delta AuditLine
\where
hookInstalled = ztrue \\
hookInstalled' = zfalse \\
spooled' = spooled \\
sealed' = sealed \\
committed' = committed \\
writerExists' = writerExists \\
rosterKnown' = rosterKnown \\
lostUnwarned' = lostUnwarned
\end{schema}

\subsection{SessionPurgeRefuse}

Models the unsealed-lines guard in \field{purgeOneTombstoned}
(internal/session/store.go:534--627): with lines still pending and no
\texttt{--force}, purge refuses outright and changes nothing. The guard
itself is \texttt{if !force \&\& (unsealed > 0 || probeFailed)}
(store.go:634--669); the caller that records the refusal is
store.go:527--530, the \field{didRefuse} path.

\begin{schema}{SessionPurgeRefuse}
\Xi AuditLine
\where
spooled > 0 \lor sealed > 0
\end{schema}

\subsection{SessionPurgeForce}

Models the same guard's \texttt{--force} branch: purge proceeds, but
leaves a \emph{flagged} tombstone recording the unsealed-lines condition
(docs/audit-seal.md \S Two zones, lines 96--105: an ``unsealed-lines
flag... set either when it was forced past pending lines or''). A flagged
tombstone is still a durable record a later \field{VacuumCrossCheck} pass
can find (audit\_vacuum.go:53--56, the tombstone branch), so
\field{rosterKnown} does \emph{not} drop to \field{zfalse} the way it does
under plain \field{SessionEnd} --- purge trades a live roster entry for an
equally discoverable tombstone, never for silence.

\begin{schema}{SessionPurgeForce}
\Delta AuditLine
\where
spooled' = spooled \\
sealed' = sealed \\
committed' = committed \\
writerExists' = writerExists \\
hookInstalled' = hookInstalled \\
rosterKnown' = rosterKnown \\
lostUnwarned' = lostUnwarned
\end{schema}

% ===================================================================
\section{Safety Property}
\label{sec:safety}
% ===================================================================

\textbf{AuditDurability}: no state reachable from \field{Init} by any
sequence of the operations in \S\ref{sec:ops} has
\field{lostUnwarned = ztrue}.

Unlike spec-mission-lifecycle.tex's \field{AbandonSafety}, this is not
checked as a schema invariant (\S\ref{sec:state} explains why it must not
be one). It is checked the way a genuine counter-example search is always
checked: by running the actual operation set through ProB and inspecting
the states it reaches.

\subsection{The current-design counter-example}
\label{sec:counterexample}

\texttt{/z-spec:check} passes on this document. \texttt{/z-spec:test}
against \field{Init} and the eight operations of \S\ref{sec:ops} explores
17 states and 78 transitions, is deadlock-free, and raises no assertion
violation --- expected, since \field{lostUnwarned} carries no invariant
for probcli to check (\S\ref{sec:state}). Reachability is instead
confirmed by a goal-directed search, \fpath{probcli ... -model\_check
-goal "lostUnwarned=ztrue" -check\_goal}, which finds
\field{lostUnwarned = ztrue} after exploring 7 states and 35 transitions,
by the shortest possible route --- no actor need do anything but remove
the writer checkout while lines are still spooled. This is probcli's own
trace, verbatim:

\begin{verbatim}
*** TRACE (length=3):
 1: SETUP_CONSTANTS(totalLines=2)
 2: INITIALISATION(committed=0,hookInstalled=ztrue,lostUnwarned=zfalse,
                    rosterKnown=ztrue,sealed=0,spooled=2,writerExists=ztrue)
 3: WorktreeRemove(lostUnwarned=ztrue,spooled=0,writerExists=zfalse)

==> lostUnwarned = ztrue
TRUE
\end{verbatim}

Two lines are destroyed in a single step, with no seal, no commit, and no
warning at the time of destruction --- \field{CommitInSiblingWorktree}'s
warning, if it ever fires, fires later and only if some other checkout
happens to commit afterward (\S\ref{sec:ops}). This is the ethos-015a loss
state the mission brief names. Recorded tool run: see
\S\ref{sec:probrun}.

% ===================================================================
\section{Alternative Designs}
\label{sec:alternatives}
% ===================================================================

Three amendments are evaluated against the same property. Each is a
self-contained variant --- its own state schema and its own operations ---
run independently of this document's own checked namespace, for the same
reason spec-mission-lifecycle.tex's two decoupled demonstrations were: to
get an isolated ProB run per variant without redeclaring \field{AuditLine}
and its operations a second (or fourth) time in one namespace, and to keep
each variant's own invariant from silently absorbing the very transition
being tested (\S\ref{sec:state}). The Z fragments below are shown verbatim
for the record; they were type-checked and model-checked as independent
\fpath{.tex} scratch files, and \S\ref{sec:verdict} reports the results.

\subsection{Design 1: per-machine spool outside any checkout}
\label{sec:design1}

The ethos-ss2x direction: move the live spool out of
\fpath{<repo>/.punt-labs/local/} entirely, into a per-machine store (a
SQLite database keyed by repo identity and session, not by checkout path).
A line is then never inside any one checkout's working tree, so
\fpath{git worktree remove} cannot reach it structurally, regardless of
which checkout is removed or when. Sealing becomes a property of
\emph{a} commit happening anywhere on the machine for the right repo, not
of a specific checkout surviving.

\begin{verbatim}
AuditLineD1
  spooled, sealed, committed : NAT
  writerExists, hookInstalled, rosterKnown, lostUnwarned : ZBOOL
where
  spooled + sealed + committed <= totalLines
  -- no writerExists-linked clause: the pool is not the writer
  -- checkout's to lose.

InitD1
  AuditLineD1'
where
  spooled' = totalLines
  sealed' = 0
  committed' = 0
  writerExists' = ztrue
  hookInstalled' = ztrue
  rosterKnown' = ztrue
  lostUnwarned' = zfalse

CommitAnywhereD1
  Delta AuditLineD1
where
  hookInstalled = ztrue
  spooled' = 0
  sealed' = 0
  committed' = committed + spooled + sealed
  writerExists' = writerExists
  hookInstalled' = hookInstalled
  rosterKnown' = rosterKnown
  lostUnwarned' = lostUnwarned

CommitAnywhereNoHookD1
  Xi AuditLineD1
where
  hookInstalled = zfalse
  -- the NoHook counterpart carries over from the current design
  -- (S:ops) unchanged in shape: still a pure query, now retyped
  -- against AuditLineD1.

WorktreeRemoveD1
  Xi AuditLineD1
where
  writerExists = ztrue
  -- pure frame: removing A checkout has nothing left to take from
  -- the per-machine spool.

-- EthosDisable, SessionEnd, SessionPurgeRefuse, SessionPurgeForce carry
-- over from the current design unchanged; none of them ever assigns
-- lostUnwarned' := ztrue, and CommitAnywhereD1/WorktreeRemoveD1 above
-- are the only two operations that could.
\end{verbatim}

\subsection{Design 2: worktree-aware seal}
\label{sec:design2}

\field{listRepoSessions} (audit\_seal.go:68) is extended to enumerate every
worktree's session directories via \fpath{git worktree list}, not only the
committing checkout's own. A commit in a sibling worktree can then seal
\emph{and commit} the writer checkout's pending lines too, provided the
writer checkout still exists to be read from.

\begin{verbatim}
AuditLineD2
  spooled, sealed, committed : NAT
  writerExists, hookInstalled, rosterKnown, lostUnwarned : ZBOOL
where
  spooled + sealed + committed <= totalLines
  writerExists = zfalse ==> (spooled = 0 and sealed = 0)
  -- identical to the current design: the pool still lives inside the
  -- writer checkout's working tree.

InitD2   -- identical in shape to Init (S:init)

CommitInWritingCheckoutD2   -- identical to CommitInWritingCheckout (S:opcommitwriting)

CommitInWritingCheckoutNoHookD2   -- identical to CommitInWritingCheckoutNoHook

CommitInSiblingWorktreeD2
  Delta AuditLineD2
where
  writerExists = ztrue
  -- worktree-aware seal can only enumerate a worktree that still
  -- exists on disk to list and read.
  spooled' = 0
  sealed' = 0
  committed' = committed + spooled + sealed
  writerExists' = writerExists
  hookInstalled' = hookInstalled
  rosterKnown' = rosterKnown
  lostUnwarned' = lostUnwarned

WorktreeRemoveD2   -- IDENTICAL to WorktreeRemove (S:opworktreeremove):
                   -- worktree-aware seal changes what a COMMIT does, not
                   -- what `git worktree remove` does, and no hook fires
                   -- on removal in this design either.

-- EthosDisable, SessionEnd, SessionPurgeRefuse, SessionPurgeForce
-- carry over unchanged.
\end{verbatim}

\subsection{Design 3: sealing SessionEnd}
\label{sec:design3}

\field{HandleSessionEnd} gains an unsealed-line probe mirroring
\field{purgeOneTombstoned}'s (store.go:534--627), but seals rather than
refuses: \field{HandleSessionEnd} already documents its cleanup as
best-effort with no retry path (session\_end.go:41--45), and a hook with
no interactive caller to retry against cannot usefully \emph{refuse} the
way a CLI command can return nonzero and let an operator re-run it.
Sealing converts a pending line from ``in the gitignored live file'' to
``in a tracked-but-uncommitted chunk file'' --- which changes what
\fpath{git worktree remove} sees: git's own dirty-worktree check refuses a
plain (non-\texttt{--force}) removal when a tracked path has uncommitted
content (\S\ref{sec:nongoals}), so a sealed-but-uncommitted chunk is
something git itself can now block on --- but only a block, not a
guarantee: \field{WorktreeRemoveD3}'s guard,
\field{force? = ztrue} $\lor$ \field{sealed = 0}, is a disjunction, so
an operator who passes \texttt{--force} keeps the operation enabled
regardless of \field{sealed}'s value, by the guard's own shape and with
no need to search for a witness state. This is the same shape of
override \field{SessionPurgeForce} already offers on the purge path
(\S\ref{sec:ops}); \S\ref{sec:verdict} returns to the difference.

\begin{verbatim}
AuditLineD3   -- identical fields and invariant to AuditLineD2/AuditLine

InitD3   -- identical in shape to Init (S:init)

CommitInWritingCheckoutD3   -- identical to CommitInWritingCheckout

CommitInWritingCheckoutNoHookD3   -- identical to CommitInWritingCheckoutNoHook

CommitInSiblingWorktreeD3   -- identical to the current design's
                            -- (pure Xi, S:ops): Design 3 does not touch
                            -- what a sibling commit does.

SessionEndD3
  Delta AuditLineD3
where
  spooled' = 0
  sealed' = sealed + spooled
  committed' = committed
  writerExists' = writerExists
  hookInstalled' = hookInstalled
  rosterKnown' = zfalse
  lostUnwarned' = lostUnwarned

WorktreeRemoveD3
  Delta AuditLineD3
  force? : ZBOOL
where
  writerExists = ztrue
  force? = ztrue or sealed = 0
  -- git's own refusal: a plain removal is blocked whenever a tracked,
  -- uncommitted chunk is sitting in the checkout. Only a still-spooled
  -- (gitignored, untracked) line is invisible to this guard.
  writerExists' = zfalse
  spooled' = 0
  sealed' = 0
  committed' = committed
  hookInstalled' = hookInstalled
  rosterKnown' = rosterKnown
  (spooled > 0 or sealed > 0) ==> lostUnwarned' = ztrue
  (spooled = 0 and sealed = 0) ==> lostUnwarned' = lostUnwarned

-- EthosDisable, SessionPurgeRefuse, SessionPurgeForce carry over
-- unchanged.
\end{verbatim}

% ===================================================================
\section{Tool Run}
\label{sec:probrun}
% ===================================================================

\begin{center}
\begin{tabularx}{\textwidth}{@{}lccX@{}}
\toprule
\textbf{Model} & \textbf{States} & \textbf{Transitions} &
  \textbf{Result} \\
\midrule
Current design (\S\ref{sec:state}--\ref{sec:ops}), all 8 operations,
  full exploration &
  17 & 78 &
  deadlock-free, no assertion violated (none declared, \S\ref{sec:state}) \\
Current design, goal search \field{lostUnwarned=ztrue} &
  6 & 37 &
  \field{lostUnwarned = ztrue} reachable via \field{WorktreeRemove} from
  \field{Init} in one step (\S\ref{sec:counterexample}) \\
Design 1 (\S\ref{sec:design1}): per-machine spool, full exploration &
  9 & 42 &
  \emph{No counter example found. ALL states visited} ---
  \field{lostUnwarned = ztrue} \textbf{unreachable} \\
Design 2 (\S\ref{sec:design2}): worktree-aware seal, full exploration &
  17 & 70 &
  deadlock-free \\
Design 2, goal search \field{lostUnwarned=ztrue} &
  4 & 23 &
  \textbf{still reachable} --- via \field{WorktreeRemoveD2} from
  \field{Init} in one step, before any
  \field{CommitInSiblingWorktreeD2} \\
Design 3 (\S\ref{sec:design3}): sealing \field{SessionEnd}, full
  exploration & 17 & 84 & deadlock-free \\
Design 3, goal search \field{lostUnwarned=ztrue} $\land$
  \field{rosterKnown=ztrue} &
  7 & 40 &
  \textbf{still reachable} --- via \field{WorktreeRemoveD3(force?=zfalse)}
  from \field{Init} in one step, with \field{SessionEndD3} never having
  fired: a spooled (gitignored) line is invisible to git's
  dirty-worktree guard, which only inspects \field{sealed} \\
\bottomrule
\end{tabularx}
\end{center}

fuzz's \texttt{/z-spec:check} reports 0 errors on this document and on
each of the three scratch variants. Every run in this table was executed
directly with \fpath{probcli}: the four full-exploration rows via
\fpath{-model\_check}, the four reachability rows via
\fpath{-model\_check -goal "<predicate>" -check\_goal}, which halts and
prints a trace the first time it finds a state satisfying the goal (a
smaller state/transition count than the corresponding full-exploration
row is expected --- the goal search stops early).

Full-exploration counts are deterministic --- re-running any of the four
full-exploration rows reproduces the same state and transition count
every time, on any probcli build, because exhaustive search visits every
reachable state regardless of the order it visits them in. Goal-search
counts are not: \fpath{-check\_goal} stops at the \emph{first} state it
happens to visit that satisfies the predicate, and re-running the exact
same goal search against the exact same unmodified file has been
observed to report a different intermediate state/transition count from
one run to the next (the current-design goal search above, run twice
across this document's two review rounds, reported 7/35 and then 6/37).
The reachability verdict --- \textbf{TRUE}, with a concrete trace ---
is what the goal search establishes; its state/transition count is a
report of how much of the space that particular run happened to touch
before stopping, not a fact about the specification. All runs were
against \fpath{.tex} files --- \texttt{probcli} silently reports 0 states
and 0 transitions on any other extension (spec-mission-lifecycle.tex \S Tool
run makes the same point; it is worth repeating because it is exactly the
kind of mistake that produces a false-clean report with no error message
at all).

% ===================================================================
\section{Verdict}
\label{sec:verdict}
% ===================================================================

Only Design~1 satisfies \field{AuditDurability} unconditionally: moving the
live spool out of any one checkout's working tree is what makes
\fpath{git worktree remove} structurally unable to reach it, in every
ordering ProB explored. This is the strongest of the three amendments and
the one this specification recommends as the fix bead ethos-ss2x should
implement --- a per-machine store, not a per-checkout hook.

Design~2 (worktree-aware seal) and Design~3 (sealing \field{SessionEnd})
each close a real ordering but not the one the mission brief's counter-example
uses. Design~2 only helps when a sibling checkout happens to commit
\emph{before} the writer checkout is removed; a \fpath{git worktree remove}
run with no other commit pending anywhere still destroys the pool exactly
as today, because worktree-aware seal changes what a \emph{commit} does,
and nothing changes what \fpath{git worktree remove} does. Design~3 only
helps when \field{HandleSessionEnd} actually ran and sealed before the
checkout was removed; a checkout whose session crashed, or was torn down
by a script that never went through a clean session end, presents the
same unsealed, gitignored pool that git's own dirty-worktree check cannot
see (confirmed: \field{lostUnwarned = ztrue} $\land$
\field{rosterKnown = ztrue} is reachable in one step from \field{Init},
\S\ref{sec:probrun}). And even in the ordering Design~3 does cover, its
refusal is only git's default one: an operator who passes
\texttt{--force} to \fpath{git worktree remove} defeats it exactly as
plainly as \field{SessionPurgeForce} lets \texttt{--force} defeat
\field{SessionPurgeRefuse} (\S\ref{sec:ops}) --- the difference is that
purge's force path still leaves a flagged tombstone, and a forced
\fpath{git worktree remove} leaves nothing. Neither is wrong to build --- each narrows the exposure window at low
cost, and Design~3 in particular converts a silent loss into a git-level
refusal for the ordering it does cover, which is exactly the kind of
promotion \field{SessionPurgeForce} (\S\ref{sec:ops}) already gets right
for the purge path. But neither is a substitute for removing the
checkout-dependency itself. The recommendation to ethos-ss2x is Design~1
as the fix, with Design~3 worth adding afterward as cheap defense in depth
for the ordering it does close; Design~2 is not worth the
\field{listRepoSessions} complexity it would add once Design~1 has already
made the underlying loss impossible.

% ===================================================================
\section{Traceability}
\label{sec:trace}
% ===================================================================

\begin{center}
\begin{tabularx}{\textwidth}{@{}llX@{}}
\toprule
\textbf{Z construct} & \textbf{Go / doc source} & \textbf{Note} \\
\midrule
\field{AuditLine} zone split &
  docs/audit-seal.md \S Two zones, lines 65--73 &
  live (gitignored) vs.\ sealed (tracked) paths \\
\field{CommitInWritingCheckout} & \field{SealRepo}
  (audit\_seal.go:49--64), \field{sealDirLocked}
  (audit\_seal.go:177--220) & seal-then-commit, same commit \\
\field{CommitInSiblingWorktree} & \field{listRepoSessions}
  (audit\_seal.go:68), \field{VacuumCrossCheck}
  (audit\_vacuum.go:46--124) & seals nothing outside \field{repoRoot};
  warns post-hoc \\
\field{WorktreeRemove} & \emph{no ethos hook exists} &
  internal/resolve only reads worktree layout
  (resolve.go:757--827); confirmed absent by inspection \\
\field{SessionEnd} & \field{HandleSessionEnd}
  (session\_end.go:46--67), \field{ss.Delete} $\to$ \field{deleteFiles}
  (internal/session/store.go:278--301) & no seal, no probe, by design \\
\field{EthosDisable} & \fpath{internal/enable/enable.go}: \field{unchainHooks}
  (line 193), info-only unsealed-count message (lines 197--198) & \\
\field{SessionPurgeRefuse}, \field{SessionPurgeForce} &
  \field{purgeOneTombstoned} (internal/session/store.go:534--627) &
  refuse-or-flagged-tombstone, never silent \\
\field{AuditDurability} counter-example &
  ethos-015a & \field{lostUnwarned} reachable in one step from
  \field{Init} \\
Design 1 & bead ethos-ss2x (SQLite per-machine spool) & recommended fix \\
Design 2 & \field{listRepoSessions} worktree-aware extension &
  partial, ordering-dependent \\
Design 3 & \field{HandleSessionEnd} seal-on-end &
  partial, crash-ordering-dependent \\
\bottomrule
\end{tabularx}
\end{center}

\end{document}
